\hypertarget{chapter-1-foundations-and-compilation}{%
\section{CHAPTER 1: FOUNDATIONS AND
COMPILATION}\label{chapter-1-foundations-and-compilation}}

\hypertarget{foundations-compilation-model}{%
\section{FOUNDATIONS \& COMPILATION
MODEL}\label{foundations-compilation-model}}

\hypertarget{absolute-basics-c98}{%
\section{ABSOLUTE BASICS (C++98)}\label{absolute-basics-c98}}

\hypertarget{getting-started}{%
\subsection{Getting Started}\label{getting-started}}

\hypertarget{what-is-c}{%
\subsubsection{What is C++?}\label{what-is-c}}

C++ is a statically-typed, compiled programming language that combines
low-level memory manipulation with high-level abstractions. It's the
language of choice for performance-critical applications.

\begin{quote}
\textbf{Brain Power: Why C++?} Think of C++ as the ``Power Tool'' of
programming. Python is like a high-end digital camerapress a button, and
it does everything for you. C++ is like a professional film camera where
you manually adjust the aperture, shutter speed, and focus. Its harder
to use, but it gives you absolute control over the final result. If
youre building a rocket, a game engine, or a high-frequency trading
system, you don't want a ``press here'' tool; you want C++.
\end{quote}

\hypertarget{your-first-program-c98}{%
\subsubsection{Your First Program
(C++98)}\label{your-first-program-c98}}

\begin{lstlisting}[language={C++}]
#include <iostream>  // 1. The Preprocessor Directive

int main() {         // 2. The Entry Point
    std::cout << "Hello, World!" << std::endl; // 3. The Output
    return 0;        // 4. The Exit Status
}
\end{lstlisting}

\hypertarget{technical-decomposition}{%
\paragraph{Technical Decomposition:}\label{technical-decomposition}}

\begin{enumerate}
\def\labelenumi{\arabic{enumi}.}
\tightlist
\item
  \textbf{\passthrough{\lstinline!\#include <iostream>!}}: This tells
  the compiler to go find the code for ``Standard Input/Output'' and
  paste it right here. Without this, the computer wouldn't know what
  \passthrough{\lstinline!std::cout!} is.
\item
  \textbf{\passthrough{\lstinline!int main()!}}: Every C++ program
  starts here. The \passthrough{\lstinline!int!} means this function
  will return an integer to the Operating System when it's done.
\item
  \textbf{\passthrough{\lstinline!std::cout!}}: Think of this as a
  ``pipe'' that leads to your screen. The \passthrough{\lstinline!<<!}
  operators are ``pushing'' the string into that pipe.
\item
  \textbf{\passthrough{\lstinline!std::endl!}}: This ends the line and
  \textbf{flushes the buffer}. Flushing the buffer is like hitting
  ``Send'' on a messageit forces the computer to actually display it
  right now.
\end{enumerate}

\begin{center}\rule{0.5\linewidth}{0.5pt}\end{center}

\hypertarget{fireside-chat-the-assembly-line-of-compilation}{%
\subsubsection{Fireside Chat: The Assembly Line of
Compilation}\label{fireside-chat-the-assembly-line-of-compilation}}

Imagine you are building a custom car. You don't just ``run'' a car; you
build it in stages. C++ works exactly the same way.

\begin{longtable}[]{@{}
  >{\raggedright\arraybackslash}p{(\columnwidth - 4\tabcolsep) * \real{0.33}}
  >{\raggedright\arraybackslash}p{(\columnwidth - 4\tabcolsep) * \real{0.33}}
  >{\raggedright\arraybackslash}p{(\columnwidth - 4\tabcolsep) * \real{0.33}}@{}}
\toprule
Stage & Analogy: The Car Factory & C++ Reality \\
\midrule
\endhead
\textbf{Preprocessing} & \textbf{The Blueprint Check}: You gather all
the parts and look at the instructions. You replace shorthand like
``Standard Engine'' with the actual full engine blueprint. & The
preprocessor looks for \passthrough{\lstinline!\#!} symbols. It pastes
in headers and expands macros. The result is one giant text file. \\
\textbf{Compilation} & \textbf{The Parts Fabrication}: You take those
blueprints and forge the raw metal into actual engine parts, wheels, and
gears. These parts are now physical, but they aren't a car yet. & The
compiler translates your C++ text into \textbf{Assembly}, which is a
low-level language the CPU understands. \\
\textbf{Assembly} & \textbf{The Component Boxing}: You put those parts
into boxes and label them. ``This box is the engine,'' ``This box is the
wheel.'' & The assembler turns assembly into \textbf{Object Files
(\passthrough{\lstinline!.o!})}. These are machine code bits that
represent your specific file. \\
\textbf{Linking} & \textbf{The Final Assembly}: You take the engine from
one box, the wheels from another, and a pre-built transmission from a
library (like Bosch or Michelin), and you bolt them all together into a
drivable car. & The linker takes all your object files and pre-built
libraries (like \passthrough{\lstinline!iostream!}) and links them into
a single \textbf{Executable}. \\
\bottomrule
\end{longtable}

\begin{quote}
\textbf{There are no dumb questions\ldots{}}

\textbf{Q: Why are there so many stages? Why can't I just ``Run'' C++
like I run Python?} \textbf{A:} Because C++ is ``AOT'' (Ahead-Of-Time)
compiled. Python is interpreted (translated as it runs). By doing all
this work upfront, C++ creates a binary that is perfectly optimized for
your specific hardware. It's like the difference between buying a
tailored suit (C++) vs.~a one-size-fits-all poncho (Python).

\textbf{Q: What happens if I forget a semicolon?} \textbf{A:} The
Compiler (Stage 2) will scream at you. Its like trying to build a car
engine with a missing boltit just won't fit together.
\end{quote}

\begin{center}\rule{0.5\linewidth}{0.5pt}\end{center}

\hypertarget{professional-notes-basics-io}{%
\subsubsection{Professional Notes: Basics \&
I/O}\label{professional-notes-basics-io}}

\hypertarget{the-compilation-pipeline-deep-dive}{%
\paragraph{1. The Compilation Pipeline (Deep
Dive)}\label{the-compilation-pipeline-deep-dive}}

Compilation in C++ is a multi-stage process that transforms
human-readable source code into machine-executable binaries. 1.
\textbf{Preprocessing}: The preprocessor (\passthrough{\lstinline!cpp!})
handles directives starting with \passthrough{\lstinline!\#!}. It
includes headers (\passthrough{\lstinline!\#include!}) and expands
macros (\passthrough{\lstinline!\#define!}). 2. \textbf{Compilation}:
The compiler (\passthrough{\lstinline!g++!},
\passthrough{\lstinline!clang++!}) translates the preprocessed source
into assembly code. 3. \textbf{Assembly}: The assembler
(\passthrough{\lstinline!as!}) converts assembly into machine-readable
object files (\passthrough{\lstinline!.o!} or
\passthrough{\lstinline!.obj!}). 4. \textbf{Linking}: The linker
(\passthrough{\lstinline!ld!}) combines multiple object files and
libraries into a single executable, resolving symbols and addresses.

\hypertarget{standard-streams-and-buffered-io}{%
\paragraph{2. Standard Streams and Buffered
I/O}\label{standard-streams-and-buffered-io}}

C++ provides a robust I/O system based on streams. *
\textbf{\passthrough{\lstinline!std::cout!}}: Buffered output stream
(standard output). * \textbf{\passthrough{\lstinline!std::cin!}}:
Buffered input stream (standard input). *
\textbf{\passthrough{\lstinline!std::cerr!}}: Unbuffered output stream
(standard error). * \textbf{\passthrough{\lstinline!std::clog!}}:
Buffered output stream (logging standard error).

\textbf{Godhood Tip}: Use \passthrough{\lstinline!std::endl!} only when
you need to flush the buffer (e.g., in real-time logging). For
performance, use \passthrough{\lstinline!'\\n'!} to avoid unnecessary
flushes.

\hypertarget{stream-state-and-error-handling}{%
\paragraph{3. Stream State and Error
Handling}\label{stream-state-and-error-handling}}

Always check the state of a stream after an operation:

\begin{lstlisting}[language={C++}]
int x;
if (!(std::cin >> x)) {
    std::cin.clear(); // Clear the error flags
    std::cin.ignore(std::numeric_limits<std::streamsize>::max(), '\n'); // Discard bad input
    std::cerr << "Invalid input received!" << std::endl;
}
\end{lstlisting}

Stream states include \passthrough{\lstinline!good()!},
\passthrough{\lstinline!bad()!}, \passthrough{\lstinline!fail()!}, and
\passthrough{\lstinline!eof()!}.

\begin{center}\rule{0.5\linewidth}{0.5pt}\end{center}

\textbf{Breakdown:} - \passthrough{\lstinline!\#include <iostream>!} -
Include input/output library - \passthrough{\lstinline!std::cout!} -
Standard output stream (print to console) -
\passthrough{\lstinline!std::endl!} - End line and flush buffer -
\passthrough{\lstinline!main()!} - Entry point of program -
\passthrough{\lstinline!return 0!} - Exit code (0 = success)

\textbf{Compile and run:}

\begin{lstlisting}[language=bash]
g++ -o hello hello.cpp
./hello
\end{lstlisting}

\begin{center}\rule{0.5\linewidth}{0.5pt}\end{center}

\hypertarget{basic-types-variables}{%
\subsection{Basic Types \& Variables}\label{basic-types-variables}}

\hypertarget{fundamental-types-c98}{%
\subsubsection{Fundamental Types (C++98)}\label{fundamental-types-c98}}

\begin{lstlisting}[language={C++}]
#include <iostream>
#include <limits>

int main() {
    // Integer types
    int x = 42;                  // 32-bit integer
    short y = 10;                // 16-bit integer
    long z = 1000000;            // 32 or 64-bit integer
    long long w = 9999999999;    // 64-bit integer

    // Floating-point types
    float f = 3.14f;             // 32-bit (4 bytes)
    double d = 3.14159265;       // 64-bit (8 bytes)
    long double ld = 3.14159265359L;  // 80+ bits

    // Character and boolean types
    char c = 'A';                // Single byte
    bool b = true;               // true or false

    // Print sizes
    std::cout << "int: " << sizeof(int) << " bytes\n";
    std::cout << "double: " << sizeof(double) << " bytes\n";

    // Min/max values
    std::cout << "int max: " << std::numeric_limits<int>::max() << "\n";
    std::cout << "int min: " << std::numeric_limits<int>::min() << "\n";

    return 0;
}
\end{lstlisting}

\begin{longtable}[]{@{}
  >{\raggedright\arraybackslash}p{(\columnwidth - 0\tabcolsep) * \real{0.06}}@{}}
\toprule
\endhead
\#\#\# Professional Notes: Literals \& Keywords \\
\#\#\#\# 1. Integer and Floating Point Literals * \textbf{Decimal}:
\passthrough{\lstinline!42!} * \textbf{Octal}:
\passthrough{\lstinline!052!} (Starts with 0) * \textbf{Hexadecimal}:
\passthrough{\lstinline!0x2A!} * \textbf{Binary (C++14)}:
\passthrough{\lstinline!0b101010!} * \textbf{Suffixes}:
\passthrough{\lstinline!u!} (unsigned), \passthrough{\lstinline!l!}
(long), \passthrough{\lstinline!ll!} (long long),
\passthrough{\lstinline!f!} (float). * \textbf{Digit Separators
(C++14)}: Use single quotes to improve readability:
\passthrough{\lstinline!1'000'000!} or
\passthrough{\lstinline!0b1111'0000!}. \\
\#\#\#\# 2. String Literals and Encodings * \textbf{Raw Strings
(C++11)}:
\passthrough{\lstinline!R"(string with \\ and " without escaping)"!} *
\textbf{Encodings}: \passthrough{\lstinline!u8""!} (UTF-8),
\passthrough{\lstinline!u""!} (char16\_t), \passthrough{\lstinline!U""!}
(char32\_t), \passthrough{\lstinline!L""!} (wchar\_t). \\
\#\#\#\# 3. Core Keywords and Type Qualifiers *
\textbf{\passthrough{\lstinline!const!}}: Declares a variable as
read-only. * \textbf{\passthrough{\lstinline!volatile!}}: Tells the
compiler that the value can change outside the program's control (e.g.,
hardware registers). Prevents aggressive optimization. *
\textbf{\passthrough{\lstinline!mutable!}}: Allows a member of a
\passthrough{\lstinline!const!} object to be modified (Chapter 5). *
\textbf{\passthrough{\lstinline!explicit!}}: Prevents the compiler from
using a constructor for implicit conversions. *
\textbf{\passthrough{\lstinline!inline!}}: Suggests to the compiler to
replace function calls with the actual code to save overhead. \\
\bottomrule
\end{longtable}

\hypertarget{variable-declaration-initialization}{%
\subsubsection{Variable Declaration \&
Initialization}\label{variable-declaration-initialization}}

\begin{lstlisting}[language={C++}]
#include <iostream>

int main() {
    // C-style initialization
    int x = 5;
    float f = 3.14f;

    // Multiple variables
    int a, b, c;

    // Uninitialized (dangerous - contains garbage)
    int uninitialized;
    std::cout << uninitialized << "\n";  // Undefined behavior!

    // Constants
    const int MAX_SIZE = 100;
    // MAX_SIZE = 200;  // Error: can't modify const

    return 0;
}
\end{lstlisting}

\hypertarget{scope-lifetime-c98}{%
\subsubsection{Scope \& Lifetime (C++98)}\label{scope-lifetime-c98}}

\begin{lstlisting}[language={C++}]
#include <iostream>

int global = 100;  // Global scope - lives entire program

void function() {
    int local = 5;      // Local scope - lives only in function
    {
        int nested = 10;  // Block scope - lives only in block
        std::cout << nested << "\n";
    }
    // std::cout << nested << "\n";  // Error: nested out of scope
}

int main() {
    {
        int x = 5;
    }
    // std::cout << x << "\n";  // Error: x out of scope

    return 0;
}
\end{lstlisting}

\hypertarget{deep-dive-the-memory-model-of-variables}{%
\subsubsection{Deep Dive: The Memory Model of
Variables}\label{deep-dive-the-memory-model-of-variables}}

Understanding \emph{where} your variables live is the first step to
Godhood.

\begin{enumerate}
\def\labelenumi{\arabic{enumi}.}
\tightlist
\item
  \textbf{The Stack (Automatic Storage)}:

  \begin{itemize}
  \tightlist
  \item
    \textbf{What}: Local variables (\passthrough{\lstinline!int x!}).
  \item
    \textbf{Speed}: Extremely fast (just moving a pointer).
  \item
    \textbf{Lifetime}: Scope-based (die at
    \passthrough{\lstinline!\}!}).
  \item
    \textbf{Limit}: Small (typically 1MB-8MB). Recursion depth is
    limited by this.
  \end{itemize}
\item
  \textbf{The Heap (Dynamic Storage)}:

  \begin{itemize}
  \tightlist
  \item
    \textbf{What}: \passthrough{\lstinline!new int!},
    \passthrough{\lstinline!malloc!}.
  \item
    \textbf{Speed}: Slower (allocation requires finding free block).
  \item
    \textbf{Lifetime}: Manual (until \passthrough{\lstinline!delete!} /
    \passthrough{\lstinline!free!}).
  \item
    \textbf{Limit}: RAM size (Gigabytes).
  \end{itemize}
\item
  \textbf{Static/Global (Static Storage)}:

  \begin{itemize}
  \tightlist
  \item
    \textbf{What}: Global variables, \passthrough{\lstinline!static!}
    locals.
  \item
    \textbf{Speed}: Fast access, but initialization order is tricky.
  \item
    \textbf{Lifetime}: Program start to program end.
  \end{itemize}
\item
  \textbf{Registers}:

  \begin{itemize}
  \tightlist
  \item
    \textbf{What}: CPU internal storage.
  \item
    \textbf{Speed}: Instant (0 cycles).
  \item
    \textbf{Note}: Variables are often optimized into registers, never
    touching RAM!
  \end{itemize}
\end{enumerate}

\begin{center}\rule{0.5\linewidth}{0.5pt}\end{center}

\hypertarget{operators-control-flow}{%
\subsection{Operators \& Control Flow}\label{operators-control-flow}}

\hypertarget{arithmetic-operators-c98}{%
\subsubsection{Arithmetic Operators
(C++98)}\label{arithmetic-operators-c98}}

\begin{lstlisting}[language={C++}]
#include <iostream>

int main() {
    int a = 10, b = 3;

    std::cout << a + b << "\n";   // 13 (addition)
    std::cout << a - b << "\n";   // 7 (subtraction)
    std::cout << a * b << "\n";   // 30 (multiplication)
    std::cout << a / b << "\n";   // 3 (integer division)
    std::cout << a % b << "\n";   // 1 (modulo)

    // Compound assignment
    int x = 5;
    x += 3;   // x = 8
    x -= 2;   // x = 6
    x *= 2;   // x = 12
    x /= 3;   // x = 4

    // Increment/Decrement
    int y = 5;
    y++;      // Post-increment: 6
    ++y;      // Pre-increment: 7
    y--;      // Post-decrement: 6
    --y;      // Pre-decrement: 5

    return 0;
}
\end{lstlisting}

\hypertarget{comparison-logical-operators-c98}{%
\subsubsection{Comparison \& Logical Operators
(C++98)}\label{comparison-logical-operators-c98}}

\begin{lstlisting}[language={C++}]
#include <iostream>

int main() {
    int a = 10, b = 5;

    // Comparison (return true/false)
    std::cout << (a > b) << "\n";   // 1 (true)
    std::cout << (a < b) << "\n";   // 0 (false)
    std::cout << (a == b) << "\n";  // 0 (false)
    std::cout << (a != b) << "\n";  // 1 (true)
    std::cout << (a >= b) << "\n";  // 1 (true)
    std::cout << (a <= b) << "\n";  // 0 (false)

    // Logical operators
    bool x = true, y = false;
    std::cout << (x && y) << "\n";  // 0 (AND)
    std::cout << (x || y) << "\n";  // 1 (OR)
    std::cout << (!x) << "\n";      // 0 (NOT)

    return 0;
}
\end{lstlisting}

\hypertarget{if-else-statements-c98}{%
\subsubsection{If-Else Statements
(C++98)}\label{if-else-statements-c98}}

\begin{lstlisting}[language={C++}]
#include <iostream>

int main() {
    int score = 85;

    // Basic if-else
    if (score >= 90) {
        std::cout << "Grade: A\n";
    } else if (score >= 80) {
        std::cout << "Grade: B\n";
    } else if (score >= 70) {
        std::cout << "Grade: C\n";
    } else {
        std::cout << "Grade: F\n";
    }

    // Ternary operator
    std::string grade = (score >= 80) ? "Pass" : "Fail";
    std::cout << grade << "\n";

    return 0;
}
\end{lstlisting}

\hypertarget{loops-c98}{%
\subsubsection{Loops (C++98)}\label{loops-c98}}

\begin{lstlisting}[language={C++}]
#include <iostream>

int main() {
    // While loop
    int i = 0;
    while (i < 5) {
        std::cout << i << " ";
        i++;
    }
    std::cout << "\n";

    // Do-while loop (executes at least once)
    int j = 0;
    do {
        std::cout << j << " ";
        j++;
    } while (j < 5);
    std::cout << "\n";

    // For loop
    for (int k = 0; k < 5; k++) {
        std::cout << k << " ";
    }
    std::cout << "\n";

    // Break and continue
    for (int m = 0; m < 10; m++) {
        if (m == 3) continue;  // Skip 3
        if (m == 7) break;     // Exit at 7
        std::cout << m << " ";
    }
    std::cout << "\n";

    return 0;
}
\end{lstlisting}

\hypertarget{switch-statement-c98}{%
\subsubsection{Switch Statement (C++98)}\label{switch-statement-c98}}

\begin{lstlisting}[language={C++}]
#include <iostream>

int main() {
    int day = 3;

    switch (day) {
        case 1:
            std::cout << "Monday\n";
            break;
        case 2:
            std::cout << "Tuesday\n";
            break;
        case 3:
            std::cout << "Wednesday\n";
            break;
        default:
            std::cout << "Unknown day\n";
    }

    return 0;
}
\end{lstlisting}

\begin{longtable}[]{@{}
  >{\raggedright\arraybackslash}p{(\columnwidth - 0\tabcolsep) * \real{0.06}}@{}}
\toprule
\endhead
\#\#\# Professional Notes: Operators \& Control Flow \\
\#\#\#\# 1. Operator Precedence and Associativity C++ has a strict
hierarchy for operator evaluation. When multiple operators appear in an
expression, precedence determines the grouping. * \textbf{High
Precedence}: \passthrough{\lstinline!()!}, \passthrough{\lstinline![]!},
\passthrough{\lstinline!->!}, \passthrough{\lstinline!::!},
\passthrough{\lstinline!++!} (postfix). * \textbf{Medium Precedence}:
\passthrough{\lstinline!*!}, \passthrough{\lstinline!/!},
\passthrough{\lstinline!\%!} followed by \passthrough{\lstinline!+!},
\passthrough{\lstinline!-!}. * \textbf{Low Precedence}: Bitwise shifts
\passthrough{\lstinline!<<!}, \passthrough{\lstinline!>>!}, then
comparisons \passthrough{\lstinline!<!}, \passthrough{\lstinline!>!},
then logical \passthrough{\lstinline!\&\&!},
\passthrough{\lstinline!||!}, and finally assignment
\passthrough{\lstinline!=!}. \\
\textbf{Godhood Warning}: Never write ambiguous code like
\passthrough{\lstinline!a = b++ + ++b;!}. This is \textbf{Undefined
Behavior (UB)} because it attempts to modify \passthrough{\lstinline!b!}
twice between sequence points. \\
\#\#\#\# 2. Advanced Loop Patterns * \textbf{Range-based for (C++11)}:
Iterate directly over containers:
\passthrough{\lstinline!for (const auto\& x : vec)!}. * \textbf{Empty
Loop Body}: A semicolon or empty braces can be used for loops that
perform all work in the header:
\passthrough{\lstinline!while (*dest++ = *src++);!}. * \textbf{The
\passthrough{\lstinline!for!} Loop as a \passthrough{\lstinline!while!}
Loop}: \passthrough{\lstinline!for (; condition ;) \{\}!} is identical
to \passthrough{\lstinline!while (condition) \{\}!}. \\
\#\#\#\# 3. Flow Control Quirks *
\textbf{\passthrough{\lstinline!switch!} Fallthrough}: Without a
\passthrough{\lstinline!break!}, execution continues to the next case.
In C++17, use \passthrough{\lstinline![[fallthrough]];!} to signal
intent and silence warnings. * \textbf{\passthrough{\lstinline!goto!}
Statement}: While generally discouraged, \passthrough{\lstinline!goto!}
is acceptable for breaking out of deeply nested loops or for
error-cleanup blocks in low-level code. * \textbf{The Comma Operator}:
\passthrough{\lstinline!a, b!} evaluates \passthrough{\lstinline!a!},
discards the result, then returns \passthrough{\lstinline!b!}. Useful in
for-loop headers:
\passthrough{\lstinline!for (int i=0, j=10; i<j; ++i, --j)!}. \\
\bottomrule
\end{longtable}

\hypertarget{functions}{%
\subsection{Functions}\label{functions}}

\hypertarget{function-declaration-definition-c98}{%
\subsubsection{Function Declaration \& Definition
(C++98)}\label{function-declaration-definition-c98}}

\begin{lstlisting}[language={C++}]
#include <iostream>

// Function declaration (prototype)
int add(int a, int b);
void print_hello();

// Function definition
int add(int a, int b) {
    return a + b;
}

void print_hello() {
    std::cout << "Hello!\n";
}

int main() {
    print_hello();
    std::cout << add(5, 3) << "\n";  // 8
    return 0;
}
\end{lstlisting}

\hypertarget{parameters-return-values-c98}{%
\subsubsection{Parameters \& Return Values
(C++98)}\label{parameters-return-values-c98}}

\begin{lstlisting}[language={C++}]
#include <iostream>

// Pass by value (copy)
void increment_value(int x) {
    x++;
    std::cout << "Inside: " << x << "\n";
}

// Pass by reference (same variable)
void increment_ref(int& x) {
    x++;
    std::cout << "Inside: " << x << "\n";
}

// Pass by const reference (can't modify)
void print_value(const int& x) {
    std::cout << x << "\n";
}

// Returning by value
int get_value() {
    return 42;
}

// Returning by reference (dangerous!)
int& get_global() {
    static int x = 100;
    return x;
}

int main() {
    int a = 5;

    increment_value(a);   // Copy passed
    std::cout << a << "\n";  // Still 5

    increment_ref(a);      // Reference passed
    std::cout << a << "\n";  // Now 6

    print_value(a);        // Can't modify a

    return 0;
}
\end{lstlisting}

\hypertarget{default-parameters-c98}{%
\subsubsection{Default Parameters
(C++98)}\label{default-parameters-c98}}

\begin{lstlisting}[language={C++}]
#include <iostream>

void greet(const std::string& name = "World") {
    std::cout << "Hello, " << name << "!\n";
}

int main() {
    greet();                // Uses default: "World"
    greet("Alice");         // Uses provided: "Alice"
    return 0;
}
\end{lstlisting}

\hypertarget{function-overloading-c98}{%
\subsubsection{Function Overloading
(C++98)}\label{function-overloading-c98}}

\begin{lstlisting}[language={C++}]
#include <iostream>

// Same function name, different parameters
int add(int a, int b) {
    return a + b;
}

double add(double a, double b) {
    return a + b;
}

void print(int x) {
    std::cout << "Integer: " << x << "\n";
}

void print(double x) {
    std::cout << "Double: " << x << "\n";
}

void print(const std::string& s) {
    std::cout << "String: " << s << "\n";
}

int main() {
    std::cout << add(5, 3) << "\n";      // 8 (int version)
    std::cout << add(2.5, 3.7) << "\n";  // 6.2 (double version)

    print(42);           // Integer version
    print(3.14);         // Double version
    print("Hello");      // String version

    return 0;
}
\end{lstlisting}

\begin{center}\rule{0.5\linewidth}{0.5pt}\end{center}

\hypertarget{arrays-pointers}{%
\subsection{Arrays \& Pointers}\label{arrays-pointers}}

\hypertarget{arrays-c98}{%
\subsubsection{Arrays (C++98)}\label{arrays-c98}}

\begin{lstlisting}[language={C++}]
#include <iostream>

int main() {
    // Array declaration and initialization
    int arr[5] = {1, 2, 3, 4, 5};

    // Access elements (0-indexed)
    std::cout << arr[0] << "\n";  // 1
    std::cout << arr[4] << "\n";  // 5

    // Array size (local arrays only)
    int size = 5;
    for (int i = 0; i < size; i++) {
        std::cout << arr[i] << " ";
    }
    std::cout << "\n";

    // 2D array
    int matrix[3][3] = {
        {1, 2, 3},
        {4, 5, 6},
        {7, 8, 9}
    };

    std::cout << matrix[1][2] << "\n";  // 6

    return 0;
}
\end{lstlisting}

\hypertarget{pointers-c98}{%
\subsubsection{Pointers (C++98)}\label{pointers-c98}}

\begin{lstlisting}[language={C++}]
#include <iostream>

int main() {
    int x = 42;

    // Create a pointer
    int* ptr = &x;  // & = address-of operator

    // Dereference pointer
    std::cout << *ptr << "\n";   // 42 (dereference with *)

    // Modify through pointer
    *ptr = 100;
    std::cout << x << "\n";      // 100

    // Pointer arithmetic
    int arr[5] = {10, 20, 30, 40, 50};
    int* p = arr;           // Array decays to pointer

    std::cout << p[0] << "\n";     // 10
    std::cout << *(p+1) << "\n";   // 20
    std::cout << p[2] << "\n";     // 30

    // Null pointer
    int* null_ptr = nullptr;  // or NULL in C++98

    // Pointer to pointer
    int** pp = &ptr;
    std::cout << **pp << "\n";  // 100

    return 0;
}
\end{lstlisting}

\hypertarget{dynamic-memory-c98}{%
\subsubsection{Dynamic Memory (C++98)}\label{dynamic-memory-c98}}

\begin{lstlisting}[language={C++}]
#include <iostream>

int main() {
    // Allocate single value on heap
    int* ptr = new int;
    *ptr = 42;
    std::cout << *ptr << "\n";
    delete ptr;          // Must deallocate
    ptr = nullptr;       // Good practice

    // Allocate array on heap
    int* arr = new int[10];
    for (int i = 0; i < 10; i++) {
        arr[i] = i * i;
    }
    delete[] arr;        // Note the [] for arrays

    // Forgetting to delete = memory leak
    int* leaked = new int(100);
    // delete leaked;  // Forgot this!

    return 0;
}
\end{lstlisting}

\hypertarget{the-c-compilation-execution-model}{%
\section{THE C++ COMPILATION \& EXECUTION
MODEL}\label{the-c-compilation-execution-model}}

To truly understand C++, you must understand how your code transforms
from text to a running process. This section demystifies the ``black
box'' of the compiler.

\hypertarget{the-build-pipeline-from-source-to-binary}{%
\subsubsection{1.5.1 The Build Pipeline: From Source to
Binary}\label{the-build-pipeline-from-source-to-binary}}

The ``compilation'' process actually consists of four distinct stages:

\begin{enumerate}
\def\labelenumi{\arabic{enumi}.}
\tightlist
\item
  \textbf{Preprocessing}: Text substitution.

  \begin{itemize}
  \tightlist
  \item
    Removes comments.
  \item
    Expands macros (\passthrough{\lstinline!\#define!}).
  \item
    Includes headers (\passthrough{\lstinline!\#include!}) recursively.
  \item
    Handles conditionals (\passthrough{\lstinline!\#ifdef!}).
  \item
    \emph{Output}: A single ``Translation Unit'' (pure C++ source code).
  \end{itemize}
\item
  \textbf{Compilation}: Syntax to Assembly.

  \begin{itemize}
  \tightlist
  \item
    \textbf{Lexical Analysis}: Tokenizes source (e.g.,
    \passthrough{\lstinline!int!}, \passthrough{\lstinline!main!},
    \passthrough{\lstinline!\{!}).
  \item
    \textbf{Parsing}: Builds Abstract Syntax Tree (AST).
  \item
    \textbf{Semantic Analysis}: Type checking, overload resolution.
  \item
    \textbf{Optimization}: Dead code elimination, loop unrolling,
    inlining (O1, O2, O3).
  \item
    \textbf{Code Generation}: Generates assembly code for the target
    architecture (x86\_64, ARM64).
  \item
    \emph{Output}: Assembly file (\passthrough{\lstinline!.s!} or
    \passthrough{\lstinline!.asm!}).
  \end{itemize}
\item
  \textbf{Assembly}: Assembly to Machine Code.

  \begin{itemize}
  \tightlist
  \item
    Translates mnemonics (\passthrough{\lstinline!MOV!},
    \passthrough{\lstinline!ADD!}) to opcodes
    (\passthrough{\lstinline!0x89!}, \passthrough{\lstinline!0x01!}).
  \item
    \emph{Output}: Object file (\passthrough{\lstinline!.o!} or
    \passthrough{\lstinline!.obj!}). Contains machine code but with
    \emph{unresolved symbols}.
  \end{itemize}
\item
  \textbf{Linking}: Combining Object Files.

  \begin{itemize}
  \tightlist
  \item
    Combines multiple \passthrough{\lstinline!.o!} files and static
    libraries
    (\passthrough{\lstinline!.a!}/\passthrough{\lstinline!.lib!}).
  \item
    Resolves symbols: Matches function calls in one TU to definitions in
    another.
  \item
    Relocates addresses: Adjusts internal pointers.
  \item
    \emph{Output}: Executable (\passthrough{\lstinline!.exe!} or
    ELF/Mach-O) or Shared Library
    (\passthrough{\lstinline!.so!}/\passthrough{\lstinline!.dll!}).
  \end{itemize}
\end{enumerate}

\hypertarget{translation-units-tu-linkage}{%
\subsubsection{1.5.2 Translation Units (TU) \&
Linkage}\label{translation-units-tu-linkage}}

A \textbf{Translation Unit} is the input to the compiler after
preprocessing. It is the fundamental unit of compilation.

\hypertarget{linkage-types}{%
\paragraph{Linkage Types}\label{linkage-types}}

Linkage determines if a name (variable/function) is visible to other
TUs.

\begin{enumerate}
\def\labelenumi{\arabic{enumi}.}
\tightlist
\item
  \textbf{No Linkage}:

  \begin{itemize}
  \tightlist
  \item
    Visible only within the current scope (block).
  \item
    Example: Local variables.
  \end{itemize}
\item
  \textbf{Internal Linkage}:

  \begin{itemize}
  \tightlist
  \item
    Visible only within the current Translation Unit.
  \item
    Hidden from the linker.
  \item
    Examples: \passthrough{\lstinline!static!} global variables,
    \passthrough{\lstinline!static!} free functions,
    \passthrough{\lstinline!const!} globals (by default).
  \item
    \emph{Use Case}: Private helper functions.
  \end{itemize}
\item
  \textbf{External Linkage}:

  \begin{itemize}
  \tightlist
  \item
    Visible to the linker and other TUs.
  \item
    Examples: Global variables, non-static functions,
    \passthrough{\lstinline!extern!} variables.
  \item
    \emph{Danger}: Requires strict ODR compliance.
  \end{itemize}
\end{enumerate}

\begin{lstlisting}[language={C++}]
// TU1.cpp
static int internal_var = 10; // Internal Linkage
int global_var = 20;          // External Linkage

// TU2.cpp
extern int global_var;        // Declares existence of external symbol
// extern int internal_var;   // ERROR: Linker error (symbol not found)
\end{lstlisting}

\hypertarget{the-one-definition-rule-odr}{%
\subsubsection{1.5.3 The One Definition Rule
(ODR)}\label{the-one-definition-rule-odr}}

The ODR is the most important rule in C++ linking.

\begin{enumerate}
\def\labelenumi{\arabic{enumi}.}
\tightlist
\item
  \textbf{ODR for Translation Units}: A function/variable can have only
  \textbf{one definition} in a single TU. (Declarations can be
  repeated).
\item
  \textbf{ODR for Programs}: A non-inline function/variable can have
  only \textbf{one definition} in the entire program.
\item
  \textbf{ODR for Classes/Templates}: Can be defined in multiple TUs
  (via headers), but definitions must be \textbf{token-for-token
  identical}.
\end{enumerate}

\textbf{Violation Example:}

\begin{lstlisting}[language={C++}]
// header.h
int x = 10; // DEFINITION (allocates memory)

// a.cpp
#include "header.h"

// b.cpp
#include "header.h"

// Linker Error: multiple definition of `x`
// Fix: Use 'extern int x;' in header, definition in one .cpp
// Fix (C++17): Use 'inline int x = 10;'
\end{lstlisting}

\hypertarget{process-memory-layout}{%
\subsubsection{1.5.4 Process Memory
Layout}\label{process-memory-layout}}

When your C++ program runs, the OS allocates virtual memory segments:

\begin{enumerate}
\def\labelenumi{\arabic{enumi}.}
\tightlist
\item
  \textbf{Text Segment (.text)}:

  \begin{itemize}
  \tightlist
  \item
    Contains executable machine instructions.
  \item
    Read-Only (prevents self-modifying code exploits).
  \item
    Shared between multiple instances of the app.
  \end{itemize}
\item
  \textbf{Data Segment (.data)}:

  \begin{itemize}
  \tightlist
  \item
    Initialized global and static variables.
  \item
    \passthrough{\lstinline!int g = 10;!} lives here.
  \end{itemize}
\item
  \textbf{BSS Segment (.bss)}:

  \begin{itemize}
  \tightlist
  \item
    Uninitialized global/static variables.
  \item
    \passthrough{\lstinline!int g;!} lives here.
  \item
    Automatically zero-initialized by OS before
    \passthrough{\lstinline!main()!}.
  \end{itemize}
\item
  \textbf{Heap}:

  \begin{itemize}
  \tightlist
  \item
    Dynamic memory (\passthrough{\lstinline!new!},
    \passthrough{\lstinline!malloc!}).
  \item
    Grows upwards (typically).
  \item
    Managed manually or by smart pointers.
  \end{itemize}
\item
  \textbf{Stack}:

  \begin{itemize}
  \tightlist
  \item
    Local variables, function parameters, return addresses.
  \item
    Grows downwards (typically).
  \item
    LIFO (Last-In, First-Out).
  \item
    \emph{Stack Overflow}: Exceeding stack size (e.g., deep recursion).
  \end{itemize}
\end{enumerate}

\hypertarget{program-startup-before-main}{%
\subsubsection{1.5.5 Program Startup: Before
main()}\label{program-startup-before-main}}

\passthrough{\lstinline!main()!} is NOT the first thing to run.

\begin{enumerate}
\def\labelenumi{\arabic{enumi}.}
\tightlist
\item
  \textbf{OS Loader}: Loads EXE into memory.
\item
  \textbf{Entry Point (\passthrough{\lstinline!\_start!})}: Defined by C
  Runtime (CRT).
\item
  \textbf{Global Initialization}:

  \begin{itemize}
  \tightlist
  \item
    Constructors of global/static objects run \textbf{before}
    \passthrough{\lstinline!main!}.
  \item
    \emph{Static Initialization Order Fiasco}: Order between TUs is
    undefined.
  \end{itemize}
\item
  \textbf{\passthrough{\lstinline!main()!}}: Your code runs.
\item
  \textbf{Global Destruction}:

  \begin{itemize}
  \tightlist
  \item
    Destructors of global/static objects run \textbf{after}
    \passthrough{\lstinline!main!} returns.
  \item
    \passthrough{\lstinline!atexit!} handlers run.
  \end{itemize}
\end{enumerate}

\hypertarget{deep-dive-into-data-representation}{%
\subsubsection{1.5.6 Deep Dive into Data
Representation}\label{deep-dive-into-data-representation}}

To understand bugs like integer overflow and floating-point inaccuracy,
you must know how data is stored bits-by-bits.

\hypertarget{integer-representation-twos-complement}{%
\paragraph{Integer Representation (Two's
Complement)}\label{integer-representation-twos-complement}}

Most modern systems use \textbf{Two's Complement} for signed integers. *
\textbf{Positive Numbers}: Standard binary.
\passthrough{\lstinline!0000 0101!} (5). * \textbf{Negative Numbers}:
Invert all bits and add 1. * \passthrough{\lstinline!5!} =
\passthrough{\lstinline!0000 0101!} * Invert:
\passthrough{\lstinline!1111 1010!} * Add 1:
\passthrough{\lstinline!1111 1011!} (-5) * \textbf{Advantage}:
Addition/Subtraction works identically for signed/unsigned. *
\textbf{Range}: \passthrough{\lstinline!-2\^(N-1)!} to
\passthrough{\lstinline!2\^(N-1) - 1!}. (One extra negative number).

\hypertarget{floating-point-ieee-754}{%
\paragraph{Floating Point (IEEE 754)}\label{floating-point-ieee-754}}

\passthrough{\lstinline!float!} (32-bit) and
\passthrough{\lstinline!double!} (64-bit) follow IEEE 754. *
\textbf{Components}: 1. \textbf{Sign Bit}: 0 (+) or 1 (-). 2.
\textbf{Exponent}: Biased integer (allows very large/small numbers). 3.
\textbf{Mantissa (Significand)}: The precision bits (normalized to
1.xxxxx). * \textbf{Implication}: *
\passthrough{\lstinline"0.1 + 0.2 != 0.3"} because 0.1 cannot be
represented exactly in binary (infinite repeating fraction). *
\textbf{NaN (Not a Number)}: Result of 0/0 or sqrt(-1).
\passthrough{\lstinline"NaN != NaN"} is always true. *
\textbf{Infinity}: Result of 1.0/0.0.

\begin{lstlisting}[language={C++}]
#include <iostream>
#include <cmath>
#include <limits>

int main() {
    float a = 0.1f;
    float b = 0.2f;
    if (a + b == 0.3f) {
        std::cout << "Math works!\n";
    } else {
        std::cout << "Math is broken: " << (a+b) << "\n"; // Prints 0.30000001
    }

    // Correct comparison
    if (std::abs((a + b) - 0.3f) < 1e-5) {
        std::cout << "Close enough!\n";
    }
}
\end{lstlisting}

\begin{center}\rule{0.5\linewidth}{0.5pt}\end{center}

\hypertarget{professional-notes-chapter-1-getting-started-with-c}{%
\section{Professional Notes: Chapter 1: Getting started with
C++}\label{professional-notes-chapter-1-getting-started-with-c}}

Version C++98 ISO/IEC 14882:1998 1998-09-01 Standard Release Date C++03
ISO/IEC 14882:2003 2003-10-16 C++11 ISO/IEC 14882:2011 2011-09-01 C++14
ISO/IEC 14882:2014 2014-12-15 C++17 TBD C++20 TBD 2017-01-01 2020-01-01
Section 1.1: Hello World This program prints Hello World! to the
standard output stream: \#include int main() \{ std::cout
\textless\textless{} ``Hello World!'' \textless\textless{} std::endl; \}
See it live on Coliru. Analysis Let's examine each part of this code in
detail: \#include is a preprocessor directive that includes the content
of the standard C++ header le iostream. iostream is a standard library
header le that contains denitions of the standard input and output
streams. These denitions are included in the std namespace, explained
below. The standard input/output (I/O) streams provide ways for programs
to get input from and output to an external system -- usually the
terminal. int main() \{ \ldots{} \} denes a new function named main. By
convention, the main function is called upon execution of the program.
There must be only one main function in a C++ program, and it must
always return a number of the int type. Here, the int is what is called
the function's return type. The value returned by the main function is
an exit code. By convention, a program exit code of 0 or EXIT\_SUCCESS
is interpreted as success by a system that executes the program. Any
other return code is associated with an error. If no return statement is
present, the main function (and thus, the program itself) returns 0 by
default. In this example, we don't need to explicitly write return 0;.
All other functions, except those that return the void type, must
explicitly return a value according to their return type, or else must
not return at all. std::cout \textless\textless{} ``Hello World!''
\textless\textless{} std::endl; prints ``Hello World!'' to the standard
output stream: std is a namespace, and :: is the scope resolution
operator that allows look-ups for objects by name within a namespace.
There are many namespaces. Here, we use :: to show we want to use cout
from the std namespace. For more information refer to Scope Resolution
Operator - Microsoft Documentation. std::cout is the standard output
stream object, dened in iostream, and it prints to the standard output
(stdout). \textless\textless{} is, in this context, the stream insertion
operator, so called because it inserts an object into the stream object.
The standard library denes the \textless\textless{} operator to perform
data insertion for certain data types into output streams. stream
\textless\textless{} content inserts content into the stream and returns
the same, but updated stream. This allows stream insertions to be
chained: std::cout \textless\textless{} ``Foo'' \textless\textless{} "
Bar``; prints''FooBar" to the console. ``Hello World!'' is a character
string literal, or a ``text literal.'' The stream insertion operator for
character string literals is dened in le iostream. std::endl is a
special I/O stream manipulator object, also dened in le iostream.
Inserting a manipulator into a stream changes the state of the stream.
The stream manipulator std::endl does two things: rst it inserts the
end-of-line character and then it ushes the stream buer to force the
text to show up on the console. This ensures that the data inserted into
the stream actually appear on your console. (Stream data is usually
stored in a buer and then ``ushed'' in batches unless you force a ush
immediately.) An alternate method that avoids the ush is: std::cout
\textless\textless{} ``Hello World!\n''; where \n is the character
escape sequence for the newline character. The semicolon (;) noties the
compiler that a statement has ended. All C++ statements and class
denitions require an ending/terminating semicolon. Section 1.2: Comments
A comment is a way to put arbitrary text inside source code without
having the C++ compiler interpret it with any functional meaning.
Comments are used to give insight into the design or method of a
program. There are two types of comments in C++: Single-Line Comments
The double forward-slash sequence // will mark all text until a newline
as a comment: int main() \{  // This is a single-line comment. int a;
// this also is a single-line comment int i; // this is another
single-line comment \} C-Style/Block Comments The sequence /* is used to
declare the start of the comment block and the sequence \emph{/ is used
to declare the end of comment. All text between the start and end
sequences is interpreted as a comment, even if the text is otherwise
valid C++ syntax. These are sometimes called ``C-style'' comments, as
this comment syntax is inherited from C++'s predecessor language, C: int
main() \{ /} * This is a block comment. \emph{/ int a; \} In any block
comment, you can write anything you want. When the compiler encounters
the symbol }/, it terminates the block comment: int main() \{ /* A block
comment with the symbol /\emph{ Note that the compiler is not affected
by the second /} however, once the end-block-comment symbol is reached,
the comment ends. \emph{/ int a; \} The above example is valid C++ (and
C) code. However, having additional /} inside a block comment might
result in a warning on some compilers. Block comments can also start and
end within a single line. For example: void SomeFunction(/* argument 1
\emph{/ int a, /} argument 2 \emph{/ int b); Importance of Comments As
with all programming languages, comments provide several benets:
Explicit documentation of code to make it easier to read/maintain
Explanation of the purpose and functionality of code Details on the
history or reasoning behind the code Placement of copyright/licenses,
project notes, special thanks, contributor credits, etc. directly in the
source code. However, comments also have their downsides: They must be
maintained to reect any changes in the code Excessive comments tend to
make the code less readable The need for comments can be reduced by
writing clear, self-documenting code. A simple example is the use of
explanatory names for variables, functions, and types. Factoring out
logically related tasks into discrete functions goes hand-in-hand with
this. Comment markers used to disable code During development, comments
can also be used to quickly disable portions of code without deleting
it. This is often useful for testing or debugging purposes, but is not
good style for anything other than temporary edits. This is often
referred to as commenting out. Similarly, keeping old versions of a
piece of code in a comment for reference purposes is frowned upon, as it
clutters les while oering little value compared to exploring the code's
history via a versioning system. Section 1.3: The standard C++
compilation process Executable C++ program code is usually produced by a
compiler. A compiler is a program that translates code from a
programming language into another form which is (more) directly
executable for a computer. Using a compiler to translate code is called
compilation. C++ inherits the form of its compilation process from its
``parent'' language, C. Below is a list showing the four major steps of
compilation in C++: 1. The C++ preprocessor copies the contents of any
included header les into the source code le, generates macro code, and
replaces symbolic constants dened using \#dene with their values. 2. The
expanded source code le produced by the C++ preprocessor is compiled
into assembly language appropriate for the platform. 3. 4. The assembler
code generated by the compiler is assembled into appropriate object code
for the platform. The object code le generated by the assembler is
linked together with the object code les for any library functions used
to produce an executable le. Note: some compiled code is linked
together, but not to create a nal program. Usually, this ``linked'' code
can also be packaged into a format that can be used by other programs.
This ``bundle of packaged, usable code'' is what C++ programmers refer
to as a library. Many C++ compilers may also merge or un-merge certain
parts of the compilation process for ease or for additional analysis.
Many C++ programmers will use dierent tools, but all of the tools will
generally follow this generalized process when they are involved in the
production of a program. The link below extends this discussion and
provides a nice graphic to help. {[}1{]}:
http://faculty.cs.niu.edu/\textasciitilde mcmahon/CS241/Notes/compile.html
Section 1.4: Function A function is a unit of code that represents a
sequence of statements. Functions can accept arguments or values and
return a single value (or not). To use a function, a function call is
used on argument values and the use of the function call itself is
replaced with its return value. Every function has a type signature --
the types of its arguments and the type of its return type. Functions
are inspired by the concepts of the procedure and the mathematical
function. Note: C++ functions are essentially procedures and do not
follow the exact denition or rules of mathematical functions. Functions
are often meant to perform a specic task. and can be called from other
parts of a program. A function must be declared and dened before it is
called elsewhere in a program. Note: popular function denitions may be
hidden in other included les (often for convenience and reuse across
many les). This is a common use of header les. Function Declaration A
function declaration is declares the existence of a function with its
name and type signature to the compiler. The syntax is as the following:
int add2(int i); // The function is of the type (int) -\textgreater{}
(int) In the example above, the int add2(int i) function declares the
following to the compiler: The return type is int. The name of the
function is add2. The number of arguments to the function is 1: The rst
argument is of the type int. The rst argument will be referred to in the
function's contents by the name i. The argument name is optional; the
declaration for the function could also be the following: int add2(int);
// Omitting the function arguments' name is also permitted. Per the
one-denition rule, a function with a certain type signature can only be
declared or dened once in an entire C++ code base visible to the C++
compiler. In other words, functions with a specic type signature cannot
be re-dened -- they must only be dened once. Thus, the following is not
valid C++: int add2(int i); // The compiler will note that add2 is a
function (int) -\textgreater{} int int add2(int j); // As add2 already
has a definition of (int) -\textgreater{} int, the compiler // will
regard this as an error. If a function returns nothing, its return type
is written as void. If it takes no parameters, the parameter list should
be empty. void do\_something(); // The function takes no parameters, and
does not return anything. // Note that it can still affect variables it
has access to. Function Call A function can be called after it has been
declared. For example, the following program calls add2 with the value
of 2 within the function of main: \#include int add2(int i); //
Declaration of add2 // Note: add2 is still missing a DEFINITION. // Even
though it doesn't appear directly in code, // add2's definition may be
LINKED in from another object file. int main() \{ std::cout
\textless\textless{} add2(2) \textless\textless{} ``\n''; // add2(2)
will be evaluated at this point, // and the result is printed. return 0;
\} Here, add2(2) is the syntax for a function call. Function Denition A
function denition} is similar to a declaration, except it also contains
the code that is executed when the function is called within its body.
An example of a function denition for add2 might be: int add2(int i) //
Data that is passed into (int i) will be referred to by the name i \{ //
while in the function's curly brackets or ``scope.'' int j = i + 2; //
Definition of a variable j as the value of i+2. return j; // Returning
or, in essence, substitution of j for a function call to // add2. \}
Function Overloading You can create multiple functions with the same
name but dierent parameters. int add2(int i) // Code contained in this
definition will be evaluated \{ // when add2() is called with one
parameter. int j = i + 2; return j; \} int add2(int i, int j) //
However, when add2() is called with two parameters, the \{ // code from
the initial declaration will be overloaded, int k = i + j + 2 ; // and
the code in this declaration will be evaluated return k; // instead. \}
Both functions are called by the same name add2, but the actual function
that is called depends directly on the amount and type of the parameters
in the call. In most cases, the C++ compiler can compute which function
to call. In some cases, the type must be explicitly stated. Default
Parameters Default values for function parameters can only be specied in
function declarations. int multiply(int a, int b = 7); // b has default
value of 7. int multiply(int a, int b) \{ return a * b; // If multiply()
is called with one parameter, the \} // value will be multiplied by the
default, 7. In this example, multiply() can be called with one or two
parameters. If only one parameter is given, b will have default value of
7. Default arguments must be placed in the latter arguments of the
function. For example: int multiply(int a = 10, int b = 20); // This is
legal int multiply(int a = 10, int b); // This is illegal since int a is
in the former Special Function Calls - Operators There exist special
function calls in C++ which have dierent syntax than
name\_of\_function(value1, value2, value3). The most common example is
that of operators. Certain special character sequences that will be
reduced to function calls by the compiler, such as !, +, -, *, \%, and
\textless\textless{} and many more. These special characters are
normally associated with non-programming usage or are used for
aesthetics (e.g.~the + character is commonly recognized as the addition
symbol both within C++ programming as well as in elementary math). C++
handles these character sequences with a special syntax; but, in
essence, each occurrence of an operator is reduced to a function call.
For example, the following C++ expression: 3+3 is equivalent to the
following function call: operator+(3, 3) All operator function names
start with operator. While in C++'s immediate predecessor, C, operator
function names cannot be assigned dierent meanings by providing
additional denitions with dierent type signatures, in C++, this is
valid. ``Hiding'' additional function denitions under one unique
function name is referred to as operator overloading in C++, and is a
relatively common, but not universal, convention in C++. Section 1.5:
Visibility of function prototypes and declarations In C++, code must be
declared or dened before usage. For example, the following produces a
compile time error: int main() \{ foo(2); // error: foo is called, but
has not yet been declared \} void foo(int x) // this later definition is
not known in main \{ \} There are two ways to resolve this: putting
either the denition or declaration of foo() before its usage in main().
Here is one example: void foo(int x) \{\} //Declare the foo function and
body first int main() \{ foo(2); // OK: foo is completely defined
beforehand, so it can be called here. \} However it is also possible to
``forward-declare'' the function by putting only a ``prototype''
declaration before its usage and then dening the function body later:
void foo(int); // Prototype declaration of foo, seen by main // Must
specify return type, name, and argument list types int main() \{ foo(2);
// OK: foo is known, called even though its body is not yet defined \}
void foo(int x) //Must match the prototype \{ // Define body of foo here
\} The prototype must specify the return type (void), the name of the
function (foo), and the argument list variable types (int), but the
names of the arguments are NOT required. One common way to integrate
this into the organization of source les is to make a header le
containing all of the prototype declarations: // foo.h void foo(int); //
prototype declaration and then provide the full denition elsewhere: //
foo.cpp --\textgreater{} foo.o \#include ``foo.h'' // foo's prototype
declaration is ``hidden'' in here void foo(int x) \{ \} // foo's body
definition and then, once compiled, link the corresponding object le
foo.o into the compiled object le where it is used in the linking phase,
main.o: // main.cpp --\textgreater{} main.o \#include ``foo.h'' // foo's
prototype declaration is ``hidden'' in here int main() \{ foo(2); \} //
foo is valid to call because its prototype declaration was beforehand.
// the prototype and body definitions of foo are linked through the
object files An unresolved external symbol error occurs when the
function prototype and call exist, but the function body is not dened.
These can be trickier to resolve as the compiler won't report the error
until the nal linking stage, and it doesn't know which line to jump to
in the code to show the error. Section 1.6: Preprocessor The
preprocessor is an important part of the compiler. It edits the source
code, cutting some bits out, changing others, and adding other things.
In source les, we can include preprocessor directives. These directives
tells the preprocessor to perform specic actions. A directive starts
with a \# on a new line. Example: \#define ZERO 0 The rst preprocessor
directive you will meet is probably the \#include directive. What it
does is takes all of something and inserts it in your le where the
directive was. The hello world program starts with the line \#include
This line adds the functions and objects that let you use the standard
input and output. The C language, which also uses the preprocessor, does
not have as many header les as the C++ language, but in C++ you can use
all the C header les. The next important directive is probably the
\#define something something\_else directive. This tells the
preprocessor that as it goes along the le, it should replace every
occurrence of something with something\_else. It can also make things
similar to functions, but that probably counts as advanced C++. The
something\_else is not needed, but if you dene something as nothing,
then outside preprocessor directives, all occurrences of something will
vanish. This actually is useful, because of the \#if,\#else and \#ifdef
directives. The format for these would be the following: \#if
something==true //code \#else //more code \#endif \#ifdef
thing\_that\_you\_want\_to\_know\_if\_is\_defined //code \#endif These
directives insert the code that is in the true bit, and deletes the
false bits. this can be used to have bits of code that are only included
on certain operating systems, without having to rewrite the whole code.

\hypertarget{professional-notes-chapter-2-literals}{%
\section{Professional Notes: Chapter 2:
Literals}\label{professional-notes-chapter-2-literals}}

Traditionally, a literal is an expression denoting a constant whose type
and value are evident from its spelling. For example, 42 is a literal,
while x is not since one must see its declaration to know its type and
read previous lines of code to know its value. However, C++11 also added
user-dened literals, which are not literals in the traditional sense but
can be used as a shorthand for function calls. Section 2.1: this Within
a member function of a class, the keyword this is a pointer to the
instance of the class on which the function was called. this cannot be
used in a static member function. struct S \{ int x; S\& operator=(const
S\& other) \{ x = other.x; // return a reference to the object being
assigned to return \emph{this; \} \}; The type of this depends on the
cv-qualication of the member function: if X::f is const, then the type
of this within f is const X}, so this cannot be used to modify
non-static data members from within a const member function. Likewise,
this inherits volatile qualication from the function it appears in.
Version C++11 this can also be used in a brace-or-equal-initializer for
a non-static data member. struct S; struct T \{ T(const S* s); //
\ldots{} \}; struct S \{ // \ldots{} T t\{this\}; \}; this is an rvalue,
so it cannot be assigned to. Section 2.2: Integer literal An integer
literal is a primary expression of the form decimal-literal It is a
non-zero decimal digit (1, 2, 3, 4, 5, 6, 7, 8, 9), followed by zero or
more decimal digits (0, 1, 2, 3, 4, 5, 6, 7, 8, 9) int d = 42;
octal-literal It is the digit zero (0) followed by zero or more octal
digits (0, 1, 2, 3, 4, 5, 6, 7) int o = 052 hex-literal It is the
character sequence 0x or the character sequence 0X followed by one or
more hexadecimal digits (0, 1, 2, 3, 4, 5, 6, 7, 8, 9, a, A, b, B, c, C,
d, D, e, E, f, F) int x = 0x2a; int X = 0X2A; binary-literal (since
C++14) It is the character sequence 0b or the character sequence 0B
followed by one or more binary digits (0, 1) int b = 0b101010; // C++14
Integer-sux, if provided, may contain one or both of the following (if
both are provided, they may appear in any order: unsigned-sux (the
character u or the character U) unsigned int u\_1 = 42u; long-sux (the
character l or the character L) or the long-long-sux (the character
sequence ll or the character sequence LL) (since C++11) The following
variables are also initialized to the same value: unsigned long long l1
= 18446744073709550592ull; // C++11 unsigned long long l2 =
18'446'744'073'709'550'592llu; // C++14 unsigned long long l3 =
1844'6744'0737'0955'0592uLL; // C++14 unsigned long long l4 =
184467'440737'0'95505'92LLU; // C++14 Notes Letters in the integer
literals are case-insensitive: 0xDeAdBaBeU and 0XdeadBABEu represent the
same number (one exception is the long-long-sux, which is either ll or
LL, never lL or Ll) There are no negative integer literals. Expressions
such as -1 apply the unary minus operator to the value represented by
the literal, which may involve implicit type conversions. In C prior to
C99 (but not in C++), unsuxed decimal values that do not t in long int
are allowed to have the type unsigned long int. When used in a
controlling expression of \#if or \#elif, all signed integer constants
act as if they have type std::intmax\_t and all unsigned integer
constants act as if they have type std::uintmax\_t. Section 2.3: true A
keyword denoting one of the two possible values of type bool. bool ok =
true; if (!f()) \{ ok = false; goto end; \} Section 2.4: false A
keyword denoting one of the two possible values of type bool. bool ok =
true; if (!f()) \{ ok = false; goto end; \} Section 2.5: nullptr Version
C++11 A keyword denoting a null pointer constant. It can be converted to
any pointer or pointer-to-member type, yielding a null pointer of the
resulting type. Widget* p = new Widget(); delete p; p = nullptr; // set
the pointer to null after deletion Note that nullptr is not itself a
pointer. The type of nullptr is a fundamental type known as
std::nullptr\_t. void f(int* p); template void g(T* p); void
h(std::nullptr\_t p); int main() \{ f(nullptr); // ok g(nullptr); //
error h(nullptr); // ok \}

\hypertarget{professional-notes-chapter-4-floating-point-arithmetic}{%
\section{Professional Notes: Chapter 4: Floating Point
Arithmetic}\label{professional-notes-chapter-4-floating-point-arithmetic}}

Section 4.1: Floating Point Numbers are Weird The rst mistake that
nearly every single programmer makes is presuming that this code will
work as intended: float total = 0; for(float a = 0; a != 2; a += 0.01f)
\{ total += a; \} The novice programmer assumes that this will sum up
every single number in the range 0, 0.01, 0.02, 0.03, \ldots, 1.97,
1.98, 1.99, to yield the result 199the mathematically correct answer.
Two things happen that make this untrue: 1. 2. The program as written
never concludes. a never becomes equal to 2, and the loop never
terminates. If we rewrite the loop logic to check a \textless{} 2
instead, the loop terminates, but the total ends up being something
dierent from 199. On IEEE754-compliant machines, it will often sum up to
about 201 instead. The reason that this happens is that Floating Point
Numbers represent Approximations of their assigned values. The classical
example is the following computation: double a = 0.1; double b = 0.2;
double c = 0.3; if(a + b == c) //This never prints on IEEE754-compliant
machines std::cout \textless\textless{} ``This Computer is Magic!''
\textless\textless{} std::endl; else std::cout \textless\textless{}
``This Computer is pretty normal, all things considered.''
\textless\textless{} std::endl; Though what we the programmer see is
three numbers written in base10, what the compiler (and the underlying
hardware) see are binary numbers. Because 0.1, 0.2, and 0.3 require
perfect division by 10which is quite easy in a base-10 system, but
impossible in a base-2 systemthese numbers have to be stored in
imprecise formats, similar to how the number 1/3 has to be stored in the
imprecise form 0.333333333333333\ldots{} in base-10. //64-bit floats
have 53 digits of precision, including the whole-number-part. double a =
0011111110111001100110011001100110011001100110011001100110011010;
//imperfect representation of 0.1 double b =
0011111111001001100110011001100110011001100110011001100110011010;
//imperfect representation of 0.2 double c =
0011111111010011001100110011001100110011001100110011001100110011;
//imperfect representation of 0.3 double a + b =
0011111111010011001100110011001100110011001100110011001100110100; //Note
that this is not quite equal to the ``canonical'' 0.3! --- \#\#\#
Professional Notes: Build Systems \& Automation

\hypertarget{the-build-process-architectural-view}{%
\paragraph{1. The Build Process (Architectural
View)}\label{the-build-process-architectural-view}}

A build system automates the invocation of the compiler, assembler, and
linker. * \textbf{Makefile}: Uses a dependency graph to determine which
files need recompilation. Only rebuilds changed files. * \textbf{CMake}:
A meta-build system that generates native build files (Makefiles, Ninja,
VS Solutions).

\hypertarget{linker-symbols-and-name-mangling}{%
\paragraph{2. Linker Symbols and Name
Mangling}\label{linker-symbols-and-name-mangling}}

Use tools like \passthrough{\lstinline!nm!} or
\passthrough{\lstinline!objdump!} to inspect object files. Demangle
names with \passthrough{\lstinline!c++filt!}.
